\documentclass[12pt,a4paper]{article}

\usepackage[numbib]{tocbibind}
\usepackage[utf8]{inputenc}
\usepackage[T1]{fontenc}
\usepackage{amsmath,amssymb}
\usepackage{graphicx}
\usepackage{booktabs}
\usepackage{longtable}
\usepackage[margin=2.5cm]{geometry}
\usepackage{setspace}
\usepackage{hyperref}
\usepackage{natbib}
\onehalfspacing

\begin{document}

% ============================================================
% COVER PAGE
% ============================================================

\begin{titlepage}
\centering
\vspace*{2cm}

{\Large Zhejiang University}\\[0.3cm]
{\large School of Economics}\\[0.5cm]
{\large Master's Thesis Proposal}\\[2cm]

{\LARGE \textbf{Compute Relocation and Energy Use in China}}\\[0.5cm]
{\large \textit{Evidence from Eastern Data, Western Compute}}\\[3cm]

\begin{tabular}{r l}
Author:      & Mario D. Martinez \\[0.3cm]
Student ID:  & 22501254 \\[0.3cm]
Program:     & Master in Chinese Economy \\[0.3cm]
Supervisor:  & Professor Du Limin \\[0.3cm]
Date:        & 18-04-2026 \\
\end{tabular}

\vfill
\end{titlepage}

\newpage
\tableofcontents
\newpage


% ============================================================
% SECTION 1 --- RESEARCH BACKGROUND AND SIGNIFICANCE
% ============================================================
% Target length: 5--7 pages
% Purpose: establish why this research matters; ground in policy
% context; state the gap the proposal fills.

\section{Research Background and Significance}

% ------------------------------------------------------------
\subsection{Research Background}

% ............................................................
\subsubsection{AI Compute Demand and Electricity Consumption}
% PURPOSE: Establish that AI adoption creates growing, material
% electricity demand at the infrastructure level.
%
% CONTENT:
% - Global data-center energy growth trajectory
% - AI training and inference as dominant new drivers of DC load
% - IEA, Masanet et al. (2020) projections
% - Most AI-economy research studies downstream outcomes;
%   this thesis focuses on the upstream energy channel
%
% REPO SUPPORT: global_plan S0, slides p.3
% STILL NEEDED: 4--5 references on AI/DC energy demand
% DO NOT OVERCLAIM: do not cite global projections as directly
% applicable to China without qualification.

Data-center electricity is the upstream channel through which AI
adoption shows up in the power system. Bottom-up accounting placed
global data-center electricity at roughly 205 TWh in 2018, about 1\%
of world electricity, and documented a striking decoupling over the
prior decade: compute instances expanded by about 550\% while
facility electricity grew only about 6\%, held down by server,
storage, and cooling efficiency gains
\citep{masanetetal-2020-recalibratingglobaldata}. That decoupling set
the baseline against which the current AI cycle is now being measured.

The AI cycle breaks that trajectory. In the United States, data-center
electricity reached 176 TWh in 2023, about 4.4\% of national
electricity, with accelerated server energy tied explicitly to the
diffusion of AI-class accelerators
\citep{shehabietal-2024-2024unitedstates}. At the global level, the
International Energy Agency projects data-center electricity to reach
on the order of 945 TWh by 2030, roughly 3\% of world electricity,
with the United States and China together accounting for about 80\%
of the growth; the same assessment flags China as a setting where the
underlying measurement base is thin \citep{internationalenergyagency-2025-energyai}.
Projections at this scale are useful for framing rather than inference
about any single Chinese province, and they reinforce rather than
substitute for province-level measurement.

The composition of AI demand determines where its electricity lands.
Training attracted early attention, but production systems shift the
long-run energy burden toward inference, and the siting of workloads
interacts with local grid and facility characteristics to determine
how much electricity is drawn and how much carbon is emitted per unit
of compute \citep{pattersonetal-2022-carbonfootprintmachine}.
Compute-infrastructure electricity is therefore not a uniform global
quantity; it is produced and consumed at specific locations whose
energy environments differ. That is the property EDWC acts on, by
concentrating new compute capacity in western Chinese provinces, and
it is the property this proposal studies. The following subsections
describe the policy that creates the shock this thesis exploits.


% ............................................................
\subsubsection{Eastern Data, Western Compute}
% PURPOSE: Describe the EDWC policy --- launch, design, goals,
% provinces involved.
%
% CONTENT:
% - February 2022 NDRC launch: 8 national hubs, 10 clusters
% - East-to-west workload migration logic: land, climate,
%   renewables, cost
% - Four treated provinces: Ningxia, Guizhou, Gansu, Inner Mongolia
% - EDWC as a bundled infrastructure intervention
%
% REPO SUPPORT: global_plan S2.1, S3; edwc_treatment_dates.md;
% NDRC source URLs
% FULLY DRAFTABLE.

China's Eastern Data, Western Compute (EDWC) program organizes new
data-center capacity around eight national computing hubs and ten
data-center clusters, with the stated aim of relocating compute
workloads from the demand-heavy eastern coast to western provinces
that have surplus land, favorable climatic conditions for cooling,
and larger shares of renewable generation
\citep{xieetal-2024-greeningchinasdigital,zhangetal-2025-easterndatawestern}.
The design rests on a reduced-form economic logic: land and power are
cheaper in the west, ambient temperatures lower cooling energy, and
provincial grid-mix differences are expected to lower the carbon
content of electricity used by western facilities. Whether that
expectation holds in practice depends on the specific source--sink
pairing, and the sign of the resulting emissions effect is not fixed
by the policy design \citep{zhangetal-2025-decarbonizingdatacenters}.

The 2023 implementation opinions on deepening EDWC moved the program
beyond initial siting guidance toward a nationally integrated computing
network, coordinating compute-power allocation, cross-regional
scheduling, transmission arrangements, and green-electricity supply.
The document rules out, in principle, construction of new large or
ultra-large data centers outside the designated hub nodes, and sets
2025 targets of more than 60\% of newly added national compute
capacity concentrated in hub-node areas and more than 80\% green
electricity in new hub-node data centers, targets that remain policy
objectives rather than realized outcomes
\citep{nationaldevelopmentandreformcommission-2023-guanyushenrushishidongshuxisuangongchengjiakuaigoujianquanguoyitihuasuanliwangdeshishiyijian}.
These targets discipline the interpretation of EDWC: the program
bundles siting, transmission, and power-mix coordination, and any
province-level effect on electricity or emissions should be read as a
reduced-form response to that bundle rather than to data-center
construction alone.

This proposal studies the four western provinces that host the
program's westernmost hubs and clusters: Ningxia, Guizhou, Gansu, and
Inner Mongolia. Each has identifiable cluster-level go-live milestones
within the 2023--2024 window, drawn from official provincial and
national sources, which provide the staggered activation timing used
in the empirical design. A fifth EDWC hub, the Chengdu--Chongqing node
centred on Sichuan, is not included in the treated set: its
pre-treatment electricity trajectory is contaminated by the
August--December 2022 hydro-thermal drought and the associated power
rationing, which is a non-EDWC shock on the same system. Go-live
milestones in the treated set are treated as working anchors in this
proposal; final locking awaits a complete source-backed provenance
pass once province-level electricity panels are in place.

Existing assessments of EDWC are either scenario-based or simulation-based.
\citet{zhangetal-2025-easterndatawestern} project emission reductions of
roughly 16--20\% by 2030 under net-zero pathway assumptions;
\citet{zhangetal-2025-decarbonizingdatacenters} simulate annual savings
of 332--942 GWh across three migration routes, with carbon outcomes
that flip sign depending on the source--sink emission-factor gap; and
\citet{xieetal-2024-greeningchinasdigital} provide scenario CO\textsubscript{2}
estimates out to 2050 using fuzzy-set qualitative comparative analysis
across hub configurations. None of these studies reports ex post,
province-level causal estimates of how EDWC go-lives have moved
electricity consumption. That is the setting the next subsections
describe.


% ............................................................
\subsubsection{Renewable Energy and Data Center Siting in Western China}
% PURPOSE: Explain why western provinces matter --- higher
% renewable share, grid differences, PUE advantages from climate.
%
% CONTENT:
% - Provincial variation in generation mix
% - Renewable curtailment as a western-province issue
% - Climate advantages for data-center cooling
%
% REPO SUPPORT: renewable_curtailment_extension.md; global_plan S4.4
% STILL NEEDED: 3--4 references on China curtailment / grid mix
% DO NOT OVERCLAIM: do not claim EDWC reduces curtailment.

The economic case for siting compute capacity in western China rests
on three provincial features: concentrated renewable generation,
climatic conditions that lower cooling energy, and electricity prices
below those on the eastern coast. None of these features translates
automatically into clean compute power. Each depends on how the local
grid absorbs, transmits, and dispatches renewable supply, and on how
facility-level efficiency is governed by policy.

Renewable abundance in the west has historically coincided with
persistent absorption problems. Wind and solar curtailment peaked in
northwestern China around 2016, with administrative and market
barriers identified as the main drivers rather than physical resource
limits \citep{liuetal-2018-curtailmentrenewableenergy}; curtailment
then declined substantially over 2015--2018 as local demand
absorption, ultra-high-voltage transmission, and conventional
interprovincial transmission contributed different shares across
provinces \citep{zhangetal-2021-whatdrivingremarkable}. More recent
provincial curtailment levels may differ from these published
figures, and the relevant point for the EDWC setting is qualitative:
the western provinces are not uniformly renewable-rich in usable
terms, because Ningxia, Guizhou, Gansu, and Inner Mongolia differ in
transmission structure and local-absorption capacity, and those
differences shape what extra electricity demand is likely to be met
from renewable rather than thermal generation.

Interregional transmission is the vehicle through which western
renewable supply reaches distant load. Optimization work on China's
multi-region power system indicates that wind and solar are projected
to account for the majority of interregional transmission volumes, and
that the northwest's outbound lines carry especially high renewable
shares; renewable accommodation in these regions depends on
transmission compatibility, flexibility resources, and the interaction
between local consumption policy and grid physics
\citep{dengetal-2022-assessingintegrationeffect}. New compute load
placed in a western hub may therefore be served partly by local
renewable generation, partly by thermal backstop, and partly by
redirecting flows that would otherwise have been exported; the mix is
a province-level question, not a national one.

Climatic conditions in several western hubs, particularly at higher
elevations and in continental interiors, support lower power usage
effectiveness through natural cooling, and this advantage is now
formalized in national policy. The 2024 action plan for green and
low-carbon data-center development sets a PUE ceiling of 1.25 for new
large and ultra-large facilities, 1.2 at national hub-node clusters,
and requires more than 80\% green electricity in new hub-node data
centers by the end of 2025
\citep{nationaldevelopmentandreformcommissionetal-2024-shujuzhongxinluseditanfazhanzhuanxiangxingdongjihua}.

These features motivate the direction of EDWC policy but do not
identify its effects. They describe why western siting was chosen and
why the electricity and emissions consequences are expected to be
heterogeneous across provinces. They are not, and should not be read
as, evidence that EDWC itself reduces curtailment or lowers provincial
carbon intensity; those are questions the empirical design takes up.


% ............................................................
\subsubsection{Policy Targets: PUE, Compute Capacity, and Carbon Peaking}
% PURPOSE: Show the policy intersection --- compute expansion +
% green DC targets + carbon peaking create the tension this
% thesis studies.
%
% CONTENT:
% - MIIT/NDRC/NEA PUE targets
% - CAICT compute-capacity targets
% - Carbon peaking / dual carbon macro constraint
% - EDWC hub clusters achieving PUE near 1.3
%
% REPO SUPPORT: global_plan S4.4 with full PUE policy source list
% FULLY DRAFTABLE.

Chinese policy now treats compute-capacity expansion and facility-level
efficiency as joint objects of regulation, held inside the macro
constraint that CO\textsubscript{2} emissions should peak before 2030
and reach net zero by 2060. The tension is direct: more compute implies
more electricity demand and, absent a cleaner grid and more efficient
facilities, more emissions. The policy response has been to push
simultaneously on the spatial allocation of compute capacity, on
facility efficiency, and on the share of green electricity used by
data centers.

Institutional reporting supplies a baseline. By the end of 2023,
Chinese data centers held roughly 8.1 million standard racks in
service, consumed on the order of 150 TWh (about 1.6\% of national
electricity), and accounted for roughly 84 Mt of
CO\textsubscript{2} emissions. The national average power usage
effectiveness fell to 1.48 in 2023 from 1.54 in 2022, and the share
of central- and western-region facilities in the rolling national
green-data-center selections rose from 26.6\% in the first batch to
46\% by the fifth, consistent with the EDWC-driven shift in where new
capacity is placed \citep{chinaacademyofinformationandcommunicationstechnologyhoringernewarea-2024-zhongguolusesuanlifazhanyanjiubaogao2024nian}.
These numbers frame the baseline against which EDWC is being evaluated
and against which this proposal's RQ2 translation is built.

The 2024 action plan tightens the efficiency and green-power targets.
By the end of 2025, the national average PUE is to fall below 1.5;
new large and ultra-large facilities are to meet a PUE of 1.25 or
lower; facilities inside hub-node clusters are held to 1.2; more than
60\% of newly added compute capacity is to be located in hub-node
areas; and green electricity is to exceed 80\% at new hub-node data
centers \citep{nationaldevelopmentandreformcommissionetal-2024-shujuzhongxinluseditanfazhanzhuanxiangxingdongjihua}.
These are policy targets; actual compliance and the province-level
distribution of compliance are empirical questions.

The policy architecture therefore binds three moving parts together:
compute-capacity growth driven by AI adoption, facility efficiency
converging toward hub-node PUE requirements, and a grid whose
provincial carbon intensities differ sharply. EDWC is the spatial
instrument that connects them, by concentrating new compute in western
hubs where cheaper land, cooler climates, and higher renewable shares
are meant to ease the trade-off. Whether that instrument has in fact
moved provincial electricity demand and carbon emissions, and how the
resulting energy flows translate into units of compute, are the
questions this proposal takes up.


% ------------------------------------------------------------
\subsection{Research Problem}
% PURPOSE: State the specific gap in one tight paragraph.
%
% KEY CLAIM: No prior study has causally estimated provincial-level
% electricity and emissions effects of EDWC go-lives using a
% staggered event-study, nor translated those effects into
% AI-relevant compute metrics.
%
% NOTE: Zhang et al. (2025) simulates EDWC carbon effects but does
% not provide causal province-level identification. The gap is
% causal DiD + compute translation.

The research problem follows from the preceding background. EDWC has
been assessed, but mainly in prospective terms: scenario-based
reductions tied to net-zero pathways and route-by-route simulations
whose sign depends on assumed source--sink emission-factor gaps
\citep{zhangetal-2025-easterndatawestern,zhangetal-2025-decarbonizingdatacenters}.
What remains missing is an ex post, province-level causal estimate of
how the program's go-lives have actually moved electricity consumption,
with CO\textsubscript{2} as a complementary environmental outcome, and
a translation of those energy effects into units of compute that
connect to the policy questions EDWC is meant to answer. Closing the
gap requires two steps the existing literature does not combine: a
staggered event-study that recovers province-specific dynamic responses
without the bias that pooled two-way fixed-effects estimators impose
under heterogeneous treatment
\citep{sunabraham-2021-estimatingdynamictreatmenta}, and a transparent
carbon-per-compute accounting layer that embeds provincial grid carbon
intensity and facility efficiency in a rate-based metric
\citep{guptaetal-2022-chasingcarbonelusive}. This proposal takes both
steps.


% ------------------------------------------------------------
\subsection{Research Significance}
% PURPOSE: Explain both theoretical and practical significance
% in one unified subsection.
%
% THEORETICAL SIGNIFICANCE:
% - Upstream AI-energy causality (vs. downstream AI-labor focus)
% - CPC as a measurement contribution
% - Applied causal inference for digital infrastructure policy
%
% PRACTICAL / POLICY SIGNIFICANCE:
% - EDWC site selection and compute-hub placement decisions
% - Green compute procurement using CPC rankings
% - Carbon accounting for AI at the provincial level
% - Implications for dual-carbon strategy
%
% REPO SUPPORT: global_plan S12

The theoretical significance of the proposed work lies in where it
sits in the AI-economy research program. Most of that work studies
downstream outcomes such as labor markets, firm productivity, and
adoption; the upstream side, where AI adoption turns into
kilowatt-hours and tonnes of CO\textsubscript{2} at specific locations,
has received less causal attention, particularly at the sub-national
level in China. A staggered event-study of EDWC go-lives contributes
an applied case in which the estimator
\citep{sunabraham-2021-estimatingdynamictreatmenta} is matched to a
compute-infrastructure shock rather than to the place-based policies
it has most often been used to evaluate. The carbon-per-compute
translation developed here extends the rate-based accounting logic
proposed for compute at the facility and workload levels
\citep{guptaetal-2022-chasingcarbonelusive} to a province-year
setting: grid carbon intensity, facility efficiency, and AI-relevant
hardware efficiency are combined into a single, reproducible series
rather than treated as separate descriptive inputs. The contribution
is measurement-side and methodological; it is not a claim about the
magnitude of any EDWC effect.

The practical significance is narrower and tied to choices that
Chinese policy is now making. Hub-node capacity allocation, green
electricity targets, and provincial compute quotas are being set
under active implementation programs
\citep{chinaacademyofinformationandcommunicationstechnologyhoringernewarea-2024-zhongguolusesuanlifazhanyanjiubaogao2024nian},
and those programs currently rely on scenario modeling rather than
ex post province-level evidence. Credible province-year estimates of
how EDWC go-lives have moved electricity consumption, with
CO\textsubscript{2} as a complementary outcome, provide a direct input
into decisions about which hubs should absorb which share of future
capacity. A transparent CPC series, built from public components,
supports green-compute procurement and carbon reporting by buyers
and operators who need a comparable rate-based number across
provinces, rather than province-specific narratives of renewable
advantage. Because provincial carbon intensities and facility
efficiencies enter the dual-carbon tracking system through different
instruments, tying them together in a single province-year series
offers a specific and reproducible contribution to how AI-related
electricity demand is accounted for inside that system. These
contributions are the ones the design can plausibly deliver; the
proposal does not extend to broader claims about the social returns
to AI or the welfare effects of compute relocation.


% ============================================================
% SECTION 2 --- LITERATURE REVIEW
% ============================================================
% Target length: 6--8 pages
% Purpose: establish what existing research has done, what it
% has not done, and why the gap matters.
%
% MUST BE SYNTHETIC, not paper-by-paper.
%
% STRUCTURAL LOGIC (derived from the research problem):
% 2.1 -- outcome side: what do we know about how compute
%        infrastructure affects electricity and emissions?
% 2.2 -- treatment/policy side: what do we know about evaluating
%        place-based infrastructure interventions like EDWC?
% 2.3 -- mechanism/translation side: what do we know about
%        converting energy into compute-relevant carbon metrics?
% 2.4 -- gap synthesis: what is missing?

\section{Literature Review}

% ------------------------------------------------------------
\subsection{Existing Research on Energy and Emissions Effects of
Compute Infrastructure}
% ANALYTICAL FUNCTION: Establish that DC/AI energy demand is
% recognized and growing, that provincial-level effects matter,
% but that causal sub-national evidence is missing.
%
% INCLUDE:
% - DC energy demand literature (IEA, Masanet et al. 2020)
% - AI training/inference energy (Patterson et al., Strubell et al.)
% - China-specific DC energy studies
% - China provincial emissions studies (Carbon Monitor, CEADs)
% - Renewable curtailment and grid-mix literature for context
% - Zhang et al. (2025) on EDWC simulation
%
% EXCLUDE: general digital-economy / productivity papers
%
% TARGET: 10--12 references, 2.5--3 pages
%
% STRUCTURE: global DC energy growth -> AI-specific demand ->
% China provincial energy/emissions -> arrive at EDWC as the
% under-studied shock on the outcome side.

Electricity demand from compute infrastructure is now a recognized
but fast-moving outcome. Global bottom-up accounting put data-center
electricity at roughly 205~TWh in 2018, close to 1\% of world
electricity, and showed that through 2018 compute-instance growth had
been largely absorbed by efficiency improvements in servers, storage,
and facility operation
\citep{masanetetal-2020-recalibratingglobaldata}. That historical
decoupling has not survived the AI acceleration. Bottom-up estimates
for the United States place data-center electricity at 176~TWh in
2023, about 4.4\% of national electricity, with 2028 scenario ranges
of roughly 325--580~TWh driven primarily by accelerated servers for
artificial intelligence \citep{shehabietal-2024-2024unitedstates}.
Global syntheses project data-center electricity near 945~TWh by 2030
in the base case, close to 3\% of world electricity, with China and
the United States jointly accounting for about 80\% of projected
growth; the same synthesis flags unusually large measurement gaps for
China, where data centers already absorb on the order of a quarter of
global data-center electricity \citep{internationalenergyagency-2025-energyai}.
The baseline claim for the outcome side is therefore that compute
infrastructure is a material and fast-growing electricity load whose
provincial distribution in China is both policy-relevant and
imperfectly observed.

A second body of work shows that the mapping from compute to
electricity to emissions is not a fixed coefficient. Bottom-up
estimates of training-footprint energy and CO\textsubscript{2} for
large neural networks vary by factors of roughly five to ten with
hardware generation, facility efficiency, and grid assumptions, so
estimates conditioned on these choices rather than on raw compute are
what the field treats as meaningful
\citep{pattersonetal-2021-carbonemissionslarge}. Follow-up work for
machine-learning operations at scale identifies data-center location
as one of a small number of first-order mitigation levers and reports
that inference, not training, often dominates operational energy in
production systems \citep{pattersonetal-2022-carbonfootprintmachine}.
Parallel cloud-instance evidence shows that selecting a cloud region
is the single most effective operational step for reducing emissions
conditional on the workload
\citep{dodgeetal-2022-measuringcarbonintensity}. Lifecycle accounting
for a single large language model confirms that inference and
deployment can add substantially to the training footprint under
realistic usage and grid mixes
\citep{luccionietal-2023-estimatingcarbonfootprint}. Two implications
follow. First, relocating workloads across grids with different
carbon intensities changes the emissions footprint even when total
electricity use is unchanged. Second, facility-efficiency differences
can shift total electricity use on their own, independent of grid
mix.

The mechanism that connects this literature to EDWC is workload
migration between regions with different renewable penetration. A
study of flexible-load migration between the CAISO and PJM footprints
finds that existing data centers can absorb a sizeable share of
regional wind and solar curtailment and reduce carbon emissions by
113--239~ktCO\textsubscript{2}e per year, while making explicit that
the sign of the emissions effect depends on the interaction between
the source grid, the destination grid, and the temporal structure of
the shifted load \citep{zhengetal-2020-mitigatingcurtailmentcarbon}.
EDWC is, in policy form, a spatial counterpart: westward relocation
of compute workloads toward provinces with higher renewable shares
and cooler climates. The relevant point for the outcome side is not
that migration is automatically clean, but that its electricity and
emissions effects must be estimated rather than assumed.

Evidence specific to China confirms that data centers already
constitute a material emissions source with pronounced provincial
variation. A Kaya--LMDI decomposition of national data-center
CO\textsubscript{2} emissions over 2017--2021 documents a marked
increase across provinces driven primarily by expanding computing
scale, with provincial emission intensities differing substantially
because of underlying grid composition and facility efficiency
\citep{nietal-2024-co2emissionmitigationpathways}. This heterogeneity
is part of what motivates evaluating EDWC province by province
rather than through a single national coefficient.

EDWC itself has been addressed in three prospective assessments.
\citet{zhangetal-2025-easterndatawestern} frame westward relocation
as a net-zero lever and estimate a roughly 16--20\% reduction in
data-center sector emissions by 2030 under scenario analysis.
\citet{zhangetal-2025-decarbonizingdatacenters} simulate energy,
economic, and carbon effects for three major east-to-west migration
routes and report annual energy savings of about 332--942~GWh from
cooling reductions and avoided transmission losses, with carbon
outcomes that vary sharply across routes: some interregional pairs
reduce emissions substantially, while pairs whose destination grid
is more carbon-intensive than the source grid increase emissions. In
both studies the sign and magnitude of the carbon effect depend on
the source--sink emission-factor gap, and neither produces an ex post
estimate of realized province-level changes around EDWC go-live
dates.

Taken together, the outcome-side literature establishes three points
that this proposal builds on. Compute-infrastructure electricity is a
material load whose trajectory has been redefined by AI. The
compute-to-emissions mapping depends on grid, facility, and
hardware, so geographic reallocation has first-order effects that are
theoretically ambiguous in sign. And within China, both the
measurement base and the small existing EDWC literature are
prospective and decomposition-based rather than causal at the
province-year level. What this literature has not delivered is the
outcome of interest for the present proposal: an ex post, province-level
estimate of how EDWC go-lives have moved provincial electricity
consumption, with CO\textsubscript{2} as a complementary outcome.


% ------------------------------------------------------------
\subsection{Existing Research on Place-Based Infrastructure Policies
and Causal Evaluation}
% ANALYTICAL FUNCTION: Show that DiD/event-study methods have been
% productively applied to infrastructure shocks, and that EDWC is
% a natural candidate. Also justify the specific methodological
% choice (Sun--Abraham over TWFE).
%
% INCLUDE:
% - Classic place-based evaluations: Kline and Moretti (2014),
%   Greenstone et al. (2010), Busso et al. (2013)
% - China-specific infrastructure/SEZ evaluations
% - Applied energy/environment DiD papers
% - Modern staggered-DiD methodology: Sun and Abraham (2021),
%   Callaway and Sant'Anna (2021), Goodman-Bacon (2021),
%   de Chaisemartin and D'Haultfoeuille (2020)
% - Abadie et al. (2010, 2015) for SCM reference (brief)
%
% EXCLUDE: pure theory of place-based policy; pre-2018 DiD
% textbook expositions
%
% TARGET: 8--10 references, 2--2.5 pages
%
% STRUCTURE: classic place-based evaluations -> China applications
% -> modern staggered-DiD methods -> gap: no causal evaluation of
% a compute-infrastructure policy like EDWC.

Place-based infrastructure programs that activate in specific
locations at specific dates have long been treated in applied
economics as quasi-experiments, and the design template that has
emerged is directly usable for a staggered infrastructure rollout
like EDWC. Canonical evaluations in this tradition share two features
relevant here: the treatment is a geographically concentrated,
time-specific intervention, and the counterfactual is built from
units that were eligible or plausibly comparable but did not receive
treatment at the same time. The Tennessee Valley Authority study uses
authorities that were proposed but never approved by Congress as
controls \citep{klinemoretti-2014-localeconomicdevelopment}; plant-
opening studies compare winner counties with runner-up counties that
narrowly lost the siting competition
\citep{--identifyingagglomerationspillovers}; and Empowerment Zone
assessments use rejected and future applicants as comparisons
\citep{bussoetal-2013-assessingincidenceefficiency}. The shared
identification logic carries directly into a
not-yet/never-treated control group under staggered adoption.

The same design approach has an established record in
China-focused applied work, and the closest methodological parallel
to the present proposal is already available in the China
energy-policy literature. An evaluation of China's economic-zone
program combines rich firm and administrative panels with a
difference-in-differences design to estimate effects of a
place-based Chinese policy, demonstrating that quasi-experimental
evaluation of geographically targeted rollouts is standard practice
in this setting
\citep{luetal-2019-placebasedpoliciescreation}. More directly,
\citet{wangetal-2023-transregionalelectricitytransmission} treat
China's ultra-high-voltage transmission projects as a quasi-natural
experiment and apply a time-varying difference-in-differences
estimator to regional carbon emissions and carbon intensity. That
paper shares the elements that matter for EDWC: a Chinese grid-linked
infrastructure shock with staggered activation dates, panel data at a
subnational level, and an energy-environment outcome. On the
outcome-variable side, \citet{deryuginaetal-2020-longrundynamicselectricity}
apply difference-in-differences to monthly residential-electricity
consumption across more than 250 Illinois communities, with explicit
parallel-trends validation, anticipation checks, and placebo tests,
providing a template for presenting dynamic electricity-consumption
effects from monthly subnational panels. The setting is a price shock
rather than an infrastructure rollout, but the outcome variable and
the empirical toolkit are the same.

The methodological literature of the last several years sharpens how
staggered event studies should be implemented. Two-way fixed-effects
specifications with leads and lags, long treated as the default in
applied work, can deliver event-time coefficients that are
contaminated by effects from other periods and apparent pre-trends
that arise purely from treatment-effect heterogeneity
\citep{sunabraham-2021-estimatingdynamictreatmenta}. Decomposition of
the pooled two-way fixed-effects estimator shows that it is a weighted
average of all possible two-group/two-period comparisons, with
weights that can be negative when treatment effects vary over time
\citep{goodman-bacon-2021-differenceindifferencesvariationtreatmenta,
dechaisemartindhaultfoeuille-2020-twowayfixedeffects}. The constructive
response has been a family of estimators that restrict comparisons to
clean pairs, typically treated cohorts against not-yet-treated or
never-treated units, and aggregate the resulting group-time average
treatment effects into event-time paths
\citep{callawaysantanna-2021-differenceindifferencesmultipletimea}.
For a staggered rollout across a small number of treated provinces
with differing go-live dates, the implication is specific: the main
estimator should be a cohort-interacted event study in the
Sun--Abraham form, with Callaway--Sant'Anna as an alternative
aggregation under the same identifying assumptions.

These three bodies of work deliver the elements the empirical design
requires: a template for treating a geographically targeted,
time-specific shock as a quasi-experiment; a China-centred precedent
that uses the same template for a grid-linked infrastructure
intervention with energy and emissions outcomes; and a
methodological consensus on how to estimate event-time effects under
staggered adoption. What the existing applied work does not contain
is a study that applies this template to a compute-infrastructure
policy of EDWC's kind, with provincial electricity consumption as
the primary outcome and CO\textsubscript{2} as a complementary
outcome, and that treats the province-level go-lives as a staggered
shock rather than a single pooled intervention.


% ------------------------------------------------------------
\subsection{Existing Research on Energy--Compute Accounting and
Carbon Metrics}
% ANALYTICAL FUNCTION: Show that translating energy/emissions into
% compute-relevant units is an emerging need, and that no existing
% framework provides a provincial-level CPC index for China.
%
% INCLUDE:
% - Carbon footprint of computing (Patterson et al. 2022,
%   Dodge et al. 2022, Wu et al. 2022)
% - PUE methodology and standards (Uptime Institute, ASHRAE)
% - Compute-efficiency benchmarking (Green500, MLPerf Power)
% - Emission-factor / carbon-intensity methodology (IPCC, CEADs)
% - China-specific: CAICT green computing reports, RMI (2024)
%
% EXCLUDE: hardware architecture papers without energy/carbon
% content; AI model papers without energy measurement
%
% TARGET: 5--7 references, 1.5--2 pages
%
% STRUCTURE: carbon footprint of AI -> PUE as site-efficiency
% metric -> compute-efficiency benchmarks -> emission factors
% and grid CI -> gap: no integrated provincial CPC framework.

Translating electricity and emissions from compute infrastructure
into a compute-relevant metric requires an accounting framework that
the computing-systems and standards literatures have now largely
converged on. The clearest formulation separates operational
emissions, which scale with electricity use and grid carbon
intensity, from embodied emissions tied to hardware manufacturing
and facility construction, and treats the footprint of a computing
system as the sum of the two expressed per unit of useful compute
rather than as a raw energy total
\citep{guptaetal-2022-chasingcarbonelusive}. The Software Carbon
Intensity specification formalizes the same logic as a rate-based
standard: operational plus embodied emissions per functional unit,
with no role for offsets as a way to lower the reported value
\citep{greensoftwarefoundation-2024-softwarecarbonintensity}. A
parallel systems literature shows that the operational term is
inherently time- and location-specific, because grid carbon intensity
varies by hour and by location with the generation mix that is
dispatched, so the same workload can produce different emissions
depending only on where and when it runs
\citep{radovanovicetal-2023-carbonawarecomputingdatacenters}. The
common structure is a compute-normalized rate whose three inputs are
grid carbon intensity, facility efficiency, and the efficiency of the
hardware itself.

Each input has its own measurement tradition. Facility efficiency is
standardly reported as Power Usage Effectiveness, defined as total
facility energy divided by the IT-equipment energy it supports, with
detailed conventions for measurement boundaries, partial-facility
variants, and reporting practice
\citep{thegreengrid-2012-puecomprehensiveexamination}. The same
practitioner literature treats PUE as an input into a broader
carbon metric rather than as a standalone environmental indicator,
because it is silent on grid mix and on hardware efficiency and
therefore cannot capture the relocation effects that motivate EDWC.

Hardware efficiency has a comparable benchmarking tradition. The
Green500 and Top500 power-measurement methodology specifies how
system power under a defined workload should be instrumented,
sampled, and reported, with graded quality levels and clear
requirements for full-run measurement, producing the audited
performance-per-watt series that now underlies most high-performance-computing
efficiency comparisons
\citep{energyefficienthighperformancecomputingworkinggroup-2015-energyefficienthigh}.
The resulting efficiency numbers are expressed in GFLOPS per watt and
provide an auditable base for the hardware-efficiency term in any
compute-normalized carbon rate; they are reported for
double-precision workloads, which implies the need for an explicit
adjustment when the target application is AI training or inference.

The grid carbon-intensity term is the input most closely tied to the
Chinese setting and the one where provincial resolution matters
most. A core method paper for China links provincial apparent energy
consumption to updated local emission factors, recalculating
provincial CO\textsubscript{2} inventories in a way that materially
changes totals relative to defaults and that now grounds the
subnational CEADs methodology
\citep{shanetal-2016-newprovincialco2}. Chinese institutional work
organizes green compute around equipment efficiency, facility
efficiency, and coordination between compute and electricity
supply, and it documents that the centre-west is now a growing share
of nationally selected green data-centre sites
\citep{chinaacademyofinformationandcommunicationstechnologyhoringernewarea-2024-zhongguolusesuanlifazhanyanjiubaogao2024nian}.
Across the standards, the systems literature, and the Chinese
institutional documents, the three inputs that a province-year
compute-normalized carbon rate requires are available as
methodologies and series, but they have not been assembled into a
single provincial Carbon-per-Compute measure for China that can be
read alongside a causal estimate of EDWC's effect on provincial
electricity consumption.


% ------------------------------------------------------------
\subsection{Literature Summary and Research Gap}
% PURPOSE: Synthesize 2.1--2.3 into a clear gap statement that
% directly motivates RQ1 and RQ2.
%
% STRUCTURE:
% - Stream 2.1 established that compute infrastructure creates
%   material energy/emissions effects, but causal sub-national
%   evidence is absent.
% - Stream 2.2 showed that event-study/DiD designs work for
%   place-based shocks, but no study has applied them to EDWC.
% - Stream 2.3 developed carbon-footprint metrics for AI, but
%   no framework provides provincial-level CPC for China
%   connecting causal energy estimates to compute units.
% - This proposal fills that gap by combining RQ1 (causal
%   province-level identification) with RQ2 (energy-to-compute
%   translation).
%
% Target: 1--1.5 pages.
% DO NOT OVERCLAIM: Zhang et al. (2025) exists. The gap is causal
% identification + CPC, not "no prior EDWC research."

The three streams reviewed above can now be read jointly. Taken
together, they establish that the outcome of interest is material,
that the identification tools for a staggered infrastructure rollout
are available, and that the measurement components for a
compute-normalized carbon rate exist. What is missing is a study that
joins these elements at the province-year level for EDWC.

On the outcome side, the evidence on energy and emissions from
compute infrastructure has two features specific to the Chinese
setting. Provincial decomposition work documents sizeable
heterogeneity in data-center emissions across regions
\citep{nietal-2024-co2emissionmitigationpathways}, and the existing
EDWC-specific assessments remain prospective, either scenario-based
\citep{zhangetal-2025-easterndatawestern} or route-simulation
\citep{zhangetal-2025-decarbonizingdatacenters}, with the sign of
the carbon effect depending on the source--sink emission-factor gap
and with workload-migration mechanisms indicating that relocation
effects are theoretically ambiguous
\citep{zhengetal-2020-mitigatingcurtailmentcarbon}. None of this
work estimates the realized, ex post effect of EDWC go-lives on
provincial electricity consumption.

On the design side, the methodological tools have matured precisely
to address the kind of rollout EDWC embodies. Event-time effects
under staggered adoption can now be estimated without the
contamination that pooled two-way fixed-effects specifications
impose under heterogeneous timing
\citep{sunabraham-2021-estimatingdynamictreatmenta,
goodman-bacon-2021-differenceindifferencesvariationtreatmenta}. The
tools have been used in China for grid-linked infrastructure and in
energy-economics settings on monthly subnational panels, but not for
a compute-infrastructure policy activating across four western
provinces over roughly two years.

On the measurement side, the accounting logic has converged, the
component methodologies are published, and Chinese institutional
documents already treat green compute as a multi-component planning
problem \citep{guptaetal-2022-chasingcarbonelusive,
chinaacademyofinformationandcommunicationstechnologyhoringernewarea-2024-zhongguolusesuanlifazhanyanjiubaogao2024nian}.
What is absent is the assembly of these components into a
province-year carbon-per-compute series for China that can be read
alongside a causal electricity estimate for the same provinces and
periods.

These three openings are not independent, and that is what justifies
pairing the proposal's two research questions rather than treating
them as separable. A causal estimate of the electricity effect
without a compute-side translation is hard to interpret against
China's stated compute-expansion goals. A compute-side translation
without a causal electricity estimate reduces to a descriptive
accounting exercise whose behavior around EDWC go-lives cannot be
attributed to the policy. The first question therefore asks how EDWC
go-lives have changed provincial electricity consumption, with
CO\textsubscript{2} as a complementary environmental outcome, using
a staggered event-study estimator rather than a pooled two-way
fixed-effects specification. The second, downstream from the first,
constructs a provincial carbon-per-compute series as a transparent
proxy, reports its levels and changes around the go-lives, and asks
what range of additional compute is compatible with the estimated
electricity effects.


% ============================================================
% SECTION 3 --- RESEARCH CONTENT AND DESIGN
% ============================================================
% Target length: 10--14 pages
% Purpose: the empirical core --- RQs, variables, data,
% identification, model, CPC framework, and feasibility.

\section{Research Content and Design}

% ------------------------------------------------------------
\subsection{Research Questions}
% RQ1: What is the causal effect of EDWC cluster go-lives on
% provincial electricity consumption and CO2 emissions?
%   - Electricity is the PRIMARY outcome.
%   - CO2 is a COMPLEMENTARY environmental outcome and input
%     into carbon intensity / CPC.
%
% RQ2: Given the provincial energy environment (CI, PUE, GF/W),
% what are the levels and changes of Carbon-per-Compute (CPC)
% before and after EDWC, and what range of additional compute is
% compatible with the estimated electricity effects?
%   - RQ2 is DOWNSTREAM from RQ1.
%
% Target: 0.5 page. FULLY DRAFTABLE.

The proposal addresses two research questions. Both concern the same
baseline treated set: the four western provinces whose EDWC hubs or
clusters went live over 2023 and 2024, namely Ningxia, Guizhou, Gansu,
and Inner Mongolia. Sichuan, the other EDWC hub province, is held out
of the treated set for the reasons introduced in \S1.1.2 and taken up
again in the data section. Both questions treat EDWC as a bundled
infrastructure package (cluster construction, grid reinforcement, and
associated renewable and transmission arrangements), and any estimated
effects are interpreted in reduced form.

\begin{itemize}
  \item \textbf{First research question: causal effects on provincial
        energy and emissions.} By how much did the arrival of an EDWC
        cluster change a province's electricity use and
        CO\textsubscript{2} emissions? Monthly provincial electricity
        consumption is the primary outcome, because electricity is the
        most direct channel through which new compute load shows up in
        the power system, and is expected to rise as new data-center
        capacity comes online. Monthly CO\textsubscript{2}, built from
        the Carbon Monitor China panel, is a complementary
        environmental outcome whose sign is theoretically ambiguous:
        load expansion pushes emissions up, while a cleaner local
        generation mix can pull them down. Each treated province is
        compared to its own pre-treatment trajectory and to provinces
        that are not yet treated or never treated within the study
        window. The design recovers the reduced-form response of the
        outcomes to the full EDWC package; it does not isolate a
        clean data-center-only effect, and it does not identify a
        causal effect on carbon intensity. The decomposition of the
        emissions response into volume and mix components is treated
        as an accounting exercise, not a separate causal claim.

  \item \textbf{Second research question: translation into
        carbon-per-compute.} Given each province's energy environment,
        how does the carbon cost of producing one unit of AI compute
        move around the cluster go-live, and how much additional
        compute is consistent with the electricity response recovered
        under the first question? Carbon-per-compute (CPC) is defined
        as provincial carbon intensity times facility power usage
        effectiveness, divided by an AI-equivalent hardware efficiency
        series, with units of tCO\textsubscript{2} per exaflop-hour;
        it is an intensive metric, not a measure of total AI compute.
        The question asks for levels and changes in CPC before and
        after each go-live, and for comparisons across provinces,
        including CAICT compute-tier peers. Results are reported as
        scenarios and uncertainty bands rather than single point
        estimates: the compute translation is run under several shares
        of the estimated electricity effect attributable to
        data-center load, and under a Monte Carlo draw over the
        parameters that enter PUE, carbon intensity, and the
        efficiency series.
\end{itemize}


% ------------------------------------------------------------
\subsection{Variables}
% PURPOSE: Define all outcome, treatment, control, and constructed
% variables used in the empirical design and CPC framework.
%
% STRUCTURE:
% - Outcome variables: monthly electricity (primary), monthly CO2
%   (complementary)
% - Treatment variable: EDWC go-live indicator and event time
% - Control variables: CDD, HDD, precipitation, holiday indicator
% - Constructed variables for RQ2: carbon intensity (CI),
%   PUE, GF/W, CPC

% --- Variable Summary Table ---
\begin{table}[htbp]
\centering
\caption{Variables, Definitions, and Data Sources}
\label{tab:variables}
\small
\begin{tabular}{p{3.6cm} p{2.2cm} p{5.0cm} p{3.5cm}}
\toprule
\textbf{Variable} & \textbf{Symbol} & \textbf{Definition} & \textbf{Source} \\
\midrule
Monthly electricity (primary)
  & $\text{Elec}_{p,t}$
  & Province-month total social electricity consumption (MWh)
  & Provincial statistical bulletins \\[0.2cm]
Monthly CO\textsubscript{2} (complementary)
  & $\text{CO}_{2,p,t}$
  & Province-month total CO\textsubscript{2} emissions (tonnes)
  & Carbon Monitor China \\[0.2cm]
EDWC treatment
  & $D_{p,t}$
  & Equal to one for treated province $p$ in months $t \geq E_p$, zero otherwise
  & Official go-live evidence \\[0.2cm]
Event time
  & $\ell$
  & $t - E_p$, months relative to go-live
  & Computed \\[0.2cm]
Climate and calendar controls
  & $\mathbf{X}_{p,t}$
  & Province-month controls absorbing weather and holiday variation
  & Standard meteorological and calendar sources \\[0.2cm]
Carbon intensity
  & $\text{CI}_{p,t}$
  & tCO\textsubscript{2}/MWh
  & Computed from CO\textsubscript{2} and electricity panels \\[0.2cm]
PUE
  & $\text{PUE}_{p,t}$
  & Site energy over IT energy
  & CAICT national averages and policy targets \\[0.2cm]
AI-equivalent efficiency
  & $\text{GF/W}_{\text{AI,CN}}(t)$
  & GFLOPS per watt
  & Green500 HPL frontier, HPL-MxP, China-specific adjustments \\[0.2cm]
Carbon-per-Compute
  & $\text{CPC}_{p,t}$
  & tCO\textsubscript{2} per exaflop-hour
  & $\text{CI}_{p,t} \times \text{PUE}_{p,t} \times 1000 / \text{GF/W}_{\text{AI,CN}}(t)$ \\
\bottomrule
\end{tabular}
\end{table}

% [Insert brief prose defining each variable group and
%  explaining construction choices]

Table~\ref{tab:variables} lists every variable used in the empirical
design and the carbon-per-compute framework. The variables fall into
four groups: outcomes, treatment, controls, and the constructed
series that feed the second research question. The panel is
province-month, with $p$ indexing province and $t$ indexing calendar
month; event time $\ell = t - E_p$ is measured relative to province
$p$'s go-live month $E_p$.

\begin{itemize}
  \item \textbf{Outcome variables.} The outcomes are the two monthly
        province-level series that the causal design targets directly.
        Monthly electricity consumption $\text{Elec}_{p,t}$ is the
        primary outcome, reported in MWh and built from official
        provincial statistical bulletins; it is the most direct
        upstream signal of new compute load in the power system.
        Monthly CO\textsubscript{2} emissions $\text{CO}_{2,p,t}$ are
        the complementary environmental outcome, taken from the
        Carbon Monitor China daily panel (31 provinces, 2019-01 to
        2025-11) aggregated to monthly totals in tonnes. Electricity
        is specified in logs to stabilize variance and ease
        percentage interpretation; CO\textsubscript{2} enters the
        accounting decomposition in levels, together with
        $\text{Elec}_{p,t}$, to construct provincial carbon intensity.

  \item \textbf{Treatment variable.} The treatment variable records
        whether a province has had an EDWC cluster go live, and if
        so, how long ago. $D_{p,t}$ is a province-month indicator
        equal to one for treated provinces in the set
        $\mathcal{G} = \{\text{Ningxia}, \text{Guizhou}, \text{Gansu},
        \text{Inner Mongolia}\}$ in months at or after their
        respective go-live, and zero otherwise. Go-live months $E_p$
        are coded from official operational milestones rather than
        policy announcements; baseline values and supporting
        evidence are given in \S3.3.3. These dates are working
        anchors, not locked, and are stress-tested in the robustness
        analysis. Sichuan is excluded from $\mathcal{G}$ for the
        reasons given in \S1.1.2 and \S3.1. Provinces that are not
        yet treated and provinces that remain untreated throughout the
        study window form the comparison pool.

  \item \textbf{Control variables.} The controls absorb province-month
        variation in electricity demand that is unrelated to EDWC.
        $\mathbf{X}_{p,t}$ denotes a vector of province-month climate
        and calendar controls, covering weather-driven load and
        within-year holiday displacement; the specific construction
        is left to the empirical chapter, once the baseline panel is
        locked. Province fixed effects and calendar-month fixed
        effects enter the specification directly and absorb
        time-invariant province heterogeneity and common seasonality,
        so $\mathbf{X}_{p,t}$ is intended to absorb residual
        variation rather than levels.

  \item \textbf{Constructed variables for the compute translation.}
        These variables do not enter the causal regressions; they
        convert the causal estimates of the first question into the
        carbon-per-compute metric used in the second. Provincial
        carbon intensity $\text{CI}_{p,t} = \text{CO}_{2,p,t} /
        \text{Elec}_{p,t}$, in tCO\textsubscript{2} per MWh, is
        computed directly from the two outcome series so that the
        intensity used in the translation is internally consistent
        with the estimated electricity and emissions responses.
        Facility power usage effectiveness $\text{PUE}_{p,t}$ is the
        ratio of site energy to IT energy, anchored to CAICT national
        averages and MIIT/NDRC/NEA targets and adjusted by province
        climate; targets are treated as targets, not as realized
        outcomes. The AI-equivalent hardware efficiency series
        $\text{GF/W}_{\text{AI,CN}}(t)$, in GFLOPS per watt, is a
        national monthly series constructed from the Green500 HPL
        frontier, the HPL-MxP mixed-precision speedup, a utilization
        realism factor, and a China-specific interconnect haircut
        reflecting the A800/H800 hardware regime. Carbon-per-compute
        is then defined as $\text{CPC}_{p,t} = \text{CI}_{p,t} \times
        \text{PUE}_{p,t} \times 1000 / \text{GF/W}_{\text{AI,CN}}(t)$,
        in tCO\textsubscript{2} per exaflop-hour; it is an intensive
        metric, not a measure of total AI compute.
\end{itemize}


% ------------------------------------------------------------
\subsection{Data Sources and Processing}
% PURPOSE: Describe each data source, acquisition method, and
% processing steps.
%
% Target: 2--2.5 pages.

% ............................................................
\subsubsection{Provincial Electricity Consumption}
% - Source: official provincial bulletins
% - Method: browser-assisted scraping pipeline
% - Current coverage: Gansu verified (71/72 months, 2020--2026);
%   other treated provinces and controls in progress
% - Acceptance rule: >=80% monthly coverage
% - Units standardized to MWh
%
% REPO SUPPORT: electricity_data_pipeline.md; src/electricity/gansu/
% DO NOT OVERCLAIM: state that coverage beyond Gansu is in progress.

The primary outcome variable, monthly province-level electricity
consumption, is built from official provincial statistical bulletins
rather than from any single consolidated national dataset, because no
public monthly province-level panel with adequate coverage of the
treated provinces is available off the shelf. Each bulletin reports
the province-wide total social electricity consumption
(\textit{quanshehui yongdianliang}) for a given month, released by
the provincial Department of Industry and Information Technology or
an equivalent provincial authority. Each province-month observation
is stored with source-level provenance at the bulletin level.

Collection proceeds province by province through a semi-automated
browser-assisted workflow that records and parses each bulletin page,
because direct HTTP scraping is blocked by the provincial portals'
anti-bot layer. Gansu is the first province for which the full
bulletin series has been collected and reconciled into a clean
province-month panel of total social electricity consumption; the
same workflow is being extended to Ningxia, Guizhou, Inner Mongolia,
and the untreated control provinces.

Every raw series is standardized to megawatt-hours, year-on-year
percentage changes reported in the bulletins are reconciled against
the levels, and duplicate or revised bulletins are resolved in favor
of the most recent official version. A province enters the main panel
only if at least 80\% of the study-window months are populated from
primary bulletins; provinces below that threshold are either carried
only in robustness checks or dropped, and mirrored reposts are
flagged separately from primary sources.

% ............................................................
\subsubsection{CO\textsubscript{2} Emissions}
% - Source: Carbon Monitor China (daily, 31 provinces, 2019--2025)
% - Citation: Liu et al. (2020); Figshare consolidated dataset
% - Processing: daily to monthly aggregation by province
%
% REPO SUPPORT: src/carbon_monitor/build_monthly_panel.py
% FULLY DRAFTABLE.

The complementary environmental outcome, monthly province-level
CO\textsubscript{2} emissions, is drawn from the Carbon Monitor China
dataset \citep{liuetal-2020-carbonmonitornearrealtimea}. Carbon
Monitor is a near-real-time, activity-based daily CO\textsubscript{2}
dataset built from sector-resolved inputs (hourly to daily
electricity generation, industrial production indices, daily mobility
data, flight tracks, and temperature-corrected commercial and
residential fuel use), extended to a sub-national Chinese panel
covering all 31 provinces from 2019 onward. Using Carbon Monitor
rather than national inventory totals keeps the emissions outcome at
the monthly, province-level frequency the event-study design
requires, and on the same timing grid as the electricity outcome. The
consolidated Carbon Monitor China file used here provides daily
province-sector records in million tonnes of CO\textsubscript{2} per
day, separately by sector (power, industry, ground transport,
aviation, residential), spanning 2019-01 to 2025-11 in the current
release. The panel covers the full pre-treatment window for every
treated province and provides at least seventeen post-treatment
months for each cohort; further data releases will extend the
post-treatment window as they become available.

Daily records are aggregated to province-month totals in tonnes by a
reproducible pipeline in the repository. The pipeline converts
reported values from millions of tonnes to tonnes, sums across all
sectors within each province-month (retaining an optional sector
filter for sectoral robustness runs), and writes a wide
province-month panel. The total-emissions series enters the causal
regressions as the complementary outcome and also enters provincial
carbon intensity $\text{CI}_{p,t} = \text{CO}_{2,p,t} /
\text{Elec}_{p,t}$, so that the intensity used in the compute
translation is internally consistent with the estimated electricity
and emissions responses. The sectoral disaggregation is not used in
the main specification: it is reserved for robustness, where the
power-sector series alone is used to check whether the estimated
response is concentrated in the channel most directly tied to
data-center load.

% ............................................................
\subsubsection{EDWC Treatment Timing}
% - Coding rule: prioritize official go-live milestones with
%   significant operational capacity
% - Working defaults: NX 2023-02, GZ 2023-09, GS 2024-06,
%   IM 2024-09
% - Robustness windows documented per province
%
% REPO SUPPORT: edwc_treatment_dates.md

Treatment timing is the hinge of the event-study design, so each
province's go-live date is coded from a documented, official
operational milestone rather than from the initial policy
announcement or groundbreaking. EDWC was announced as a national
program in February 2022, but the cluster-anchor facilities covered
here became operational at different dates, generating the staggered
adoption structure the empirical design exploits. The go-live month
$E_p$ for province $p$ is defined as the first calendar month in
which an official provincial or central source reports the
designated cluster-anchor facility as operating, using milestone
verbs such as \textit{shangxian yunying}, \textit{touyun}, or
\textit{put into use}, rather than verbs reporting construction
completion, acceptance, or planned trial operation. A hub-wide
operational state is not assumed; $E_p$ marks the first defensible
project-level operational milestone within the province's EDWC
cluster. In monthly panels the treatment month is coded as the first
day of the event month; in daily robustness specifications the exact
day is retained when the source is precise, and otherwise
month-level timing is used with adjacent windows as sensitivity
checks.

Baseline $T_0$ months are Ningxia 2023-02, Guizhou 2024-06, Gansu
2024-06, and Inner Mongolia 2024-02. Ningxia is coded to February
2023, when the integrated EDWC computing-power service platform in
Yinchuan was officially launched and reported as operating.\footnote{Source
documentation and full evidence tables for each baseline and
robustness anchor are maintained in the repository note
\texttt{01\_notes/edwc\_treatment\_dates.md}; non-academic sources
such as provincial government releases and national news outlets are
cited there rather than in running prose.} Guizhou is coded to June
2024, when the overall migration of the State Power Investment
Corporation Beijing data center to Gui'an was completed, marking the
cleanest documented operational switch-over for the Guizhou cluster.
Gansu is coded to June 2024, when the provincial government reported
the Qingyang cluster as built and operating at roughly 15{,}000
standard racks and compute scale above 10{,}000~P. Inner Mongolia is
coded to February 2024, the earliest phased operational milestone of
the China Mobile Hohhot Intelligent Computing Center in Horinger.

These anchors are working dates, not locked final treatment dates;
full province-by-province fact-checking is planned once the baseline
electricity and CO\textsubscript{2} panels are ready for event-study
work. Each province has at least one documented alternative milestone
that serves as a robustness $T_0$: for Ningxia, the 2023-07
acceptance of China Telecom Ningxia Data Center Phase I; for
Guizhou, the 2023-09 early-use evidence for the SPIC Gui'an Data
Center; for Gansu, the 2024-12 and 2025-01 scale-up reports on the
Qingyang cluster; for Inner Mongolia, the 2024-04 public formal
commissioning of the Hohhot Intelligent Computing Center and the
2024-09 commissioning of the Jiuzhou Intelligent Computing Center.
The main results are re-estimated on each of these alternative
anchors so that any sensitivity to the exact coded go-live month is
transparent.

% ............................................................
\subsubsection{Weather Controls and Other Inputs}
% - CDD, HDD, precipitation by province-month
% - Source TBD (CMA or ERA5 reanalysis)
% - CAICT compute indices for heterogeneity / donor selection
%
% REPO SUPPORT: caict_series_map.md

Beyond the outcome and treatment series, the empirical design uses
two further groups of province-month inputs: climate and calendar
controls that absorb weather-driven and calendar-driven variation in
electricity demand, and CAICT compute indices used for province
benchmarking, heterogeneity splits, and donor-pool selection for the
supplementary synthetic control on Gansu.

The controls in $\mathbf{X}_{p,t}$ are intended to absorb
province-month variation in electricity demand that is unrelated to
EDWC. They will draw on standard climate inputs (temperature and
precipitation measures suitable for monthly panels) and a
province-month holiday indicator capturing Spring Festival and other
national holidays whose timing moves across calendar months. The
specific functional forms, sources, and aggregation rules are left to
the empirical chapter, and are not committed here, because province
fixed effects and calendar-month fixed effects already absorb
time-invariant province heterogeneity and common seasonality; the
role of $\mathbf{X}_{p,t}$ is to absorb residual variation rather
than levels.

The CAICT comprehensive computing-power evaluation
(\textit{zonghe suanli zhishu}), in its 2022 and 2023 editions,
provides province-level indices on four dimensions: compute capacity
in exaflops on an FP32 reference, storage in exabytes, network
quality, and environmental indicators such as PUE and green-energy
share. These indices are not used as regressors in the causal design.
They are used for three auxiliary purposes: assigning treated and
control provinces to compute tiers for heterogeneity reporting,
selecting a donor pool of provinces with similar compute environments
for the supplementary synthetic control on Gansu, and providing
cross-sectional context for carbon-per-compute comparisons in the
second research question. CAICT also publishes a separate
computing-power development index
(\textit{suanli fazhan zhishu}) with a different weighting scheme and
a different set of sub-indicators; cross-edition comparisons in this
proposal stay within the comprehensive-compute framework and do not
mix the two indices.


% ------------------------------------------------------------
\subsection{Empirical Strategy and Model}
% PURPOSE: Present the full empirical design for both RQs,
% including the event-study model, the CPC translation framework,
% and robustness plans.
%
% Target: 4--5 pages total for this subsection.

% ............................................................
\subsubsection{Identification Strategy for RQ1}
% CONTENT:
% - EDWC as a staggered policy shock across 4 provinces
% - Sun--Abraham as the main estimator: avoids TWFE bias under
%   staggered adoption and heterogeneous effects
% - Identifying assumptions:
%   (a) conditional parallel trends
%   (b) no anticipation (pre-trend test)
% - Controls: province FE, time FE, weather, holidays
%
% REPO SUPPORT: EDWC_Gansu_EventStudy S5.1; global_plan S5.1
% FULLY DRAFTABLE.

The first research question is identified from the staggered go-live
of EDWC clusters across the four treated provinces. Ningxia, Guizhou,
Gansu, and Inner Mongolia enter treatment in four different months
between early 2023 and late 2024, while the remaining provinces are
untreated over the study window. This variation in adoption timing,
rather than a single pre-post contrast, is what the design exploits.

The main estimator is the cohort-interacted event study of
\citet{sunabraham-2021-estimatingdynamictreatmenta}. For each cohort
$E_p$ and each event time $\ell$ it estimates the cohort-average
treatment effect $\text{CATT}_{E_p, \ell}$ by comparing the cohort to
units that are not yet treated or never treated, and then aggregates
these cohort-specific effects into the event-time path reported in
the results. The choice of a staggered-robust estimator is not
discretionary. A standard two-way fixed effects regression with leads
and lags is unsafe in this kind of heterogeneous, staggered setting:
\citet{goodman-bacon-2021-differenceindifferencesvariationtreatmenta}
shows that the pooled TWFE estimator is a weighted average of all
possible two-group, two-period comparisons, including comparisons
that use already-treated cohorts as controls for later-treated
cohorts; \citet{dechaisemartindhaultfoeuille-2020-twowayfixedeffects}
shows that under heterogeneous effects these weights can be negative,
so the pooled coefficient can be misleading even in sign; and
\citet{sunabraham-2021-estimatingdynamictreatmenta} shows that
dynamic TWFE specifications produce coefficients on a given event
time that are contaminated by effects from other event times, so
apparent pre-trends can arise purely from treatment-effect
heterogeneity.

Two assumptions support a causal reading of the estimated event-time
path. The first is conditional parallel trends: absent EDWC, treated
provinces and their not-yet-treated or never-treated comparisons
would have followed the same outcome trajectory, possibly after
conditioning on the observed controls. The second is no anticipation:
the outcome in treated provinces is not affected by the EDWC go-live
in months before it occurs. Both are probed empirically. Parallel
trends is assessed with event-time pre-coefficients ($\ell < 0$),
where a flat, jointly insignificant pre-trajectory is consistent with
the assumption; no anticipation is assessed by inspecting the
coefficients in the months immediately before $\ell = 0$ and by
re-estimating the model with placebo go-live months shifted into the
pre-period. The comparison pool for each cohort is formed from
provinces that are not yet treated at that event time and provinces
that remain untreated throughout the window, with Sichuan held out
of the treated set for the reasons given in §3.1. Province fixed
effects absorb time-invariant heterogeneity in levels; calendar-month
fixed effects absorb common seasonality and aggregate shocks shared
across provinces; the climate and holiday controls defined in §3.3.4
absorb weather-driven and calendar-driven variation in electricity
demand that is not tied to EDWC.

The multi-period
\citet{callawaysantanna-2021-differenceindifferencesmultipletimea}
estimator, which identifies group-time average treatment effects
under conditional parallel trends and aggregates them into event-time
paths, is reported alongside the main specification. Consistent signs,
magnitudes, and pre-trend diagnostics across the two families are
treated as supportive evidence; systematic divergence is flagged and
discussed. Additional robustness checks on the sample, the event
window, the coded go-live month, and the treatment of never-treated
versus not-yet-treated comparisons are detailed in §3.4.4.

% ............................................................
\subsubsection{Event-Study Model Specification}
% The Sun--Abraham estimating equation:

\begin{equation}
Y_{p,t} = \alpha_p + \delta_t
+ \sum_{e \in \mathcal{E}} \sum_{\substack{\ell \in \mathcal{L} \\ \ell \neq -1}}
  \beta_{e,\ell} \, \mathbb{1}\{E_p = e\} \cdot \mathbb{1}\{t - e = \ell\}
+ \mathbf{X}_{p,t}' \boldsymbol{\gamma}
+ \varepsilon_{p,t}
\label{eq:event_study}
\end{equation}

% DEFINITIONS:
% - Y_{p,t}: monthly outcome; electricity (primary) and CO2
%   (complementary), estimated in separate runs
% - E_p: treatment cohort for province p (go-live month coded in
%   S3.3.3); untreated provinces take E_p = infinity
% - E: set of treated cohorts; L: set of event times; l = -1 omitted
% - beta_{e,l}: cohort-specific event-time CATT
% - alpha_p: province FE; delta_t: calendar-month FE
% - X_{p,t}: climate and holiday controls defined in S3.3.4
% - epsilon_{p,t}: clustered by province

Equation~\eqref{eq:event_study} is the operational form of the
cohort-interacted event study. It is estimated separately on each of
the two outcomes, following the template for event-study presentation
used in applied electricity work
\citep{deryuginaetal-2020-longrundynamicselectricity}.

The dependent variable $Y_{p,t}$ is the monthly outcome for province
$p$ in calendar month $t$. Electricity enters in logs and
CO\textsubscript{2} in levels, and the two are estimated in separate
runs so that coefficients are reported in their native scale without
forcing a common functional form. $E_p$ denotes the treatment cohort
of province $p$, that is, the go-live month coded in §3.3.3; untreated
provinces are assigned $E_p = \infty$ and never contribute to the
treatment indicators. The set of treated cohorts is
$\mathcal{E} = \{E_p : p \in \mathcal{G}\}$, with one element per
distinct go-live month; the set of event times is $\mathcal{L}$, with
$\ell = -1$ omitted as the reference period so that every reported
coefficient is read relative to the month immediately before go-live.
The product
$\mathbb{1}\{E_p = e\} \cdot \mathbb{1}\{t - e = \ell\}$ turns on only
in calendar month $t$ for a province in cohort $e$ at event time
$\ell$. $\alpha_p$ is a province fixed effect, $\delta_t$ is a
calendar-month fixed effect, and $\mathbf{X}_{p,t}$ collects the
climate and holiday controls defined in §3.3.4.

The baseline event window is $\ell \in [-12, +18]$ months. Twelve
pre-treatment months are enough to recover a full calendar year of
seasonality for every cohort before treatment; the post-window
matches the longest post-treatment horizon currently available for
the earliest-treated cohort. Event times outside this window are not
dropped but bin-pooled into $\ell \leq -13$ and $\ell \geq +19$
endpoint indicators, so that observations outside the reported window
still contribute to the fixed effects without distorting the
event-time path. Only the coefficients inside the window are reported
and interpreted; the bin indicators are carried as nuisance terms.

The cohort-specific coefficients $\beta_{e,\ell}$ estimate cohort-average
treatment effects, $\text{CATT}_{e,\ell}$, by comparing each cohort at
event time $\ell$ to units that are not yet treated or never treated,
following \citet{sunabraham-2021-estimatingdynamictreatmenta}. For
reporting, the cohort-specific coefficients are aggregated into a
single event-time path,
\[
\hat{\beta}_\ell
= \sum_{e \in \mathcal{E}} w_e \, \hat{\beta}_{e,\ell},
\]
where the weights $w_e$ are the share of treated provinces observed
at event time $\ell$ in each cohort and sum to one. This aggregation
delivers a contamination-free event-time path in the staggered
setting described in §3.4.1; without it, a pooled TWFE specification
with leads and lags would mix the effects of other cohorts at other
event times into each $\beta_\ell$ and could even generate apparent
pre-trends from treatment-effect heterogeneity alone.

Standard errors $\varepsilon_{p,t}$ are clustered by province, which
is the level at which treatment is assigned. Province clustering
accommodates arbitrary serial correlation within a province and is
the most conservative choice given the small number of treated units.
Because the number of clusters is modest, wild-cluster bootstrap
p-values are reported alongside the asymptotic standard errors for
the event-time coefficients.

Results are presented as event-time figures with point estimates and
confidence intervals at each $\ell$, together with a pre-trend test
on the joint significance of the pre-coefficients
($\ell \in [-12, -2]$) and a placebo exercise that shifts $E_p$ into
the pre-period. The imputation estimator of
\citet{borusyaketal-2024-revisitingeventstudydesigns}, which recovers
each event-time coefficient by imputing untreated potential outcomes
from the fitted fixed-effects model on untreated observations only,
is reported as an additional specification. Agreement across the
Sun--Abraham main estimator, the
\citet{callawaysantanna-2021-differenceindifferencesmultipletimea}
estimator introduced in §3.4.1, and the
\citet{borusyaketal-2024-revisitingeventstudydesigns} imputation
estimator is treated as evidence that the result is not an artifact
of any single estimator.

% ............................................................
\subsubsection{Carbon-per-Compute Translation Framework}
% PURPOSE: Present the CPC framework as the translation layer
% that operationalizes RQ2, downstream from RQ1.
%
% CPC DEFINITION:

\begin{equation}
\text{CPC}_{p,t}
= \text{CI}_{p,t} \times \text{PUE}_{p,t}
  \times \frac{1000}{\text{GF/W}_{\text{AI,China}}(t)}
\quad \left[\frac{\text{tCO}_2}{\text{EF}\cdot\text{h}}\right]
\label{eq:cpc}
\end{equation}

% GF/W CONSTRUCTION:

\begin{equation}
\text{GF/W}_{\text{AI,China}}(t)
= \underbrace{\text{GF/W}_{\text{FP64}}(t) \cdot s_{\text{mxp}}(t)
  \cdot \kappa_{\text{global}}}_{\text{GF/W}_{\text{AI,global}}(t)}
  \times \chi_{\text{CN}}(t)
\label{eq:gfw}
\end{equation}

% CONTENT:
% - CI computed from CO2 and kWh panels (DiD-CPC consistency)
% - PUE trajectories: national anchor (~1.48 in 2023), EDWC hubs
%   converging toward 1.3, seasonal climate adjustment
% - GF/W: Green500 FP64 frontier x HPL-MxP speedup x MFU (0.5)
%   x China interconnect haircut (A800/H800 epochs)
% - Translation: define theta in [0,1] as data-center share of
%   additional MWh; compute Delta EF.h scenarios
% - Monte Carlo uncertainty over s_mxp, kappa, lambda, PUE, CI
%
% REPO SUPPORT: AI_equivalent_GFLOPS_per_W.md; global_plan S6
% FULLY DRAFTABLE from repo method notes.
%
% DO NOT OVERCLAIM: CPC is a method, not a result. Call it a
% "transparent proxy" or "reproducible approximation."

The second research question is operationalized through
carbon-per-compute (CPC), defined in Equation~\eqref{eq:cpc}. CPC
multiplies the carbon content of a unit of grid electricity
($\text{CI}_{p,t}$) by the site overhead that separates IT energy
from total facility energy ($\text{PUE}_{p,t}$), and divides by the
AI-equivalent energy productivity of current hardware
($\text{GF/W}_{\text{AI,China}}(t)$). The result is
tCO\textsubscript{2} per exaflop-hour, an intensive metric that
tracks the carbon cost of running compute, not the volume of compute
\citep{guptaetal-2022-chasingcarbonelusive,
radovanovicetal-2023-carbonawarecomputingdatacenters}. CPC can fall
because the grid decarbonizes, because facilities become more
efficient, or because hardware improves; the three channels are kept
separate by construction so that each contribution is visible. The
framework is presented as a transparent, reproducible approximation
rather than as an estimator of a single true carbon cost of compute
in China. Its purpose is to translate the reduced-form electricity
and emissions responses recovered by the event study into
compute-relevant units that can be compared across provinces and
over time.

CPC takes three inputs, each built from a distinct data source:

\begin{itemize}
  \item \textbf{Carbon intensity input.} Provincial carbon intensity
        is built from the two outcome series used in the event study,
        as $\text{CI}_{p,t} = \text{CO}_{2,p,t} / \text{kWh}_{p,t}$,
        so that the intensity feeding the translation is internally
        consistent with the estimated responses
        \citep{shanetal-2016-newprovincialco2}. A second CI
        construction from the provincial power mix and standard
        emission factors is retained for scenario work, in particular
        for counterfactual paths in which the pre-EDWC mix is held
        fixed.

  \item \textbf{PUE input.} Facility power usage effectiveness
        follows the Green Grid definition of site energy over IT
        energy \citep{thegreengrid-2012-puecomprehensiveexamination}.
        $\text{PUE}_{p,t}$ is specified as a province-month
        trajectory anchored to CAICT national averages, with EDWC
        provinces converging toward the hub target of roughly 1.3 by
        2025 and uncertainty bands of $\pm 0.1$ to $\pm 0.2$
        \citep{chinaacademyofinformationandcommunicationstechnologyhoringernewarea-2024-zhongguolusesuanlifazhanyanjiubaogao2024nian}.
        Targets are treated as targets, not as realized outcomes;
        sensitivity to the PUE path is reported alongside the central
        values.

  \item \textbf{AI-equivalent efficiency series.} The hardware
        efficiency series $\text{GF/W}_{\text{AI,China}}(t)$ is
        constructed in the two layers written in
        Equation~\eqref{eq:gfw}. The global layer starts from the
        Green500 HPL FP64 frontier, measured under the standard
        methodology of the Energy Efficient High Performance Computing
        Working Group
        \citep{energyefficienthighperformancecomputingworkinggroup-2015-energyefficienthigh},
        multiplies it by the HPL-MxP mixed-precision speedup
        \citep{dongarraluszczek-2026-hplmxpbenchmarkmixedprecision}
        to move from FP64 to AI-like mixed-precision workloads, and
        applies a utilization realism factor
        $\kappa_{\text{global}} = 0.5$ consistent with reported
        large-scale model FLOPs utilization in the 40\% to 60\% range
        \citep{fernandezetal-2025-efficienthardwarescaling}. The
        China-specific layer applies an interconnect haircut
        $\chi_{\text{CN}}(t) = (1 - \lambda) + \lambda \, \beta(t)$,
        where $\beta(t)$ is the ratio of export-control-compliant to
        reference-SKU GPU bandwidth (about 0.667 for A800 in
        2022-11 to 2023-10 and about 0.444 for H800 from 2023-11
        onward) and $\lambda \approx 0.4$ is the communication-bound
        share of workload execution.\footnote{The bandwidth ratios
        and the communication-bound share are placeholder calibrations
        taken from the repository notes in
        \texttt{01\_notes/AI\_equivalent\_GFLOPS\_per\_W.md}; each
        figure will be attached to a primary source (vendor
        specification sheets for A800 and H800, and the measurement
        or benchmarking study used for $\lambda$) before the
        translation is reported in the empirical chapter.}

\end{itemize}

The translation step takes the event-study estimates
$\overline{\Delta\text{MWh}}_{p}$ and
$\overline{\Delta\text{CO}_{2}}_{p}$ and maps them into additional
compute in three complementary views. All three are accounting
relationships between the causal estimates and compute, not separate
causal identifications:

\begin{itemize}
  \item \textbf{Energy-based view.} Introduce $\theta \in [0,1]$, the
        data-center share of the additional provincial MWh, and
        report
        $\Delta\text{EF}\cdot\text{h} = \theta \,
        \overline{\Delta\text{MWh}}_{p} / e^{\text{Site}}_{p,\text{post}}$,
        where $e^{\text{Site}}_{p,\text{post}} = \text{PUE}_{p,t} /
        \text{GF/W}_{\text{AI,China}}(t)$ is the post-treatment
        site-energy intensity of compute. This view measures the
        extra compute that could have been supported by the
        estimated load response and is invariant to changes in
        carbon intensity.

  \item \textbf{Emissions-compatible view.} Report the compute
        consistent with the estimated emissions response at the
        post-treatment CPC. A net variant uses the observed
        post-treatment CI and returns the compute delivered by the
        realized change in emissions, while a gross variant holds CI
        at its pre-treatment value and isolates the compute
        attributable to load expansion net of generation-mix changes.

  \item \textbf{Emissions-offset index.} Summarize, for each
        treated province, the share of the load-induced potential
        emissions that was absorbed by a cleaner generation mix.
        Define $\Delta\text{CO}_2^{\text{volume}}_{p} =
        \overline{\Delta\text{MWh}}_{p} \cdot \text{CI}_{p,\text{pre}}$
        as the counterfactual emissions change that would have
        followed mechanically from the estimated load response had
        the pre-treatment carbon intensity stayed in place, and
        compare it to the estimated emissions response
        $\overline{\Delta\text{CO}_2}_{p}$. The index
        $\text{Offset}_{p} = 1 -
        \overline{\Delta\text{CO}_2}_{p} / \Delta\text{CO}_2^{\text{volume}}_{p}$
        measures how much of the volume-implied emissions were
        absorbed by a cleaner mix: a value of one means the full
        load response was absorbed, zero means none, and negative
        values mean the mix moved the wrong way. The index is
        informative only when $\overline{\Delta\text{MWh}}_{p}$ is
        bounded away from zero, since otherwise the denominator
        collapses; it is reported alongside, not in place of, the
        energy-based and emissions-compatible views.
\end{itemize}

The translation is propagated through a Monte Carlo over the
parameters that enter CPC and the compute scenarios:
$s_{\text{mxp}}(t)$, $\kappa_{\text{global}}$, $\lambda$,
$\text{PUE}_{p,t}$, and $\text{CI}_{p,t}$. Simple uniform or
triangular distributions are specified for each parameter,
realisations of $\text{GF/W}_{\text{AI,China}}(t)$,
$e^{\text{Site}}_{p,t}$, and $\text{CPC}_{p,t}$ are drawn and
propagated through the three translation views, and results are
reported as P10 to P90 bands with the median rather than as single
point estimates. The event-study coefficients themselves enter the
Monte Carlo through their estimated sampling distribution, so the
reported bands combine parameter uncertainty in the compute bridge
with sampling uncertainty in the causal estimates.

% ............................................................
\subsubsection{Robustness and Heterogeneity}
% PURPOSE: Concise list of planned robustness exercises.
%
% CONTENT:
% - Alternative T0 windows (+/-1--2 months)
% - Placebo provinces and placebo years
% - SCM for Gansu as a supplementary robustness exercise (brief;
%   reference Abadie et al. 2010, 2015; not a co-equal design)
% - Alternative PUE/GF/W parameters for CPC sensitivity
% - Heterogeneity by CAICT compute tier and by province
% - Exclusion of confounded provinces (Sichuan 2022)
%
% Target: 0.5--1 page. FULLY DRAFTABLE as a plan.

The event-study estimates and the CPC translation are each subjected
to a set of robustness and heterogeneity exercises. The purpose is not
to vindicate a single headline number but to show how the results move
under changes in the coded treatment date, the comparison set, the
event window, the estimator, and the parameters that enter the
compute bridge.

Sensitivity to the coded go-live month is assessed by shifting $E_p$
by $\pm 1$ and $\pm 2$ months around the baseline date, using the
calendar- and announcement-based alternatives already recorded for
each treated province in §3.3.3, and by estimating the model on two
cuts of the treated set: the full four-province set and a restricted
set that drops each treated province in turn. Placebo tests assign
fake treatment dates one and two years before each true $E_p$ and
re-estimate Equation~\eqref{eq:event_study} on the pre-treatment
sample; the placebo event-time path should be flat and statistically
indistinguishable from zero if the identifying variation is genuinely
tied to go-live timing. A spatial placebo assigns the EDWC treatment
to non-EDWC provinces in the control pool and checks that
pseudo-coefficients are centred on zero. Event-window robustness
re-estimates the model with shorter ($[-6, +12]$) and longer
($[-18, +24]$) windows where data permit, and with alternative tail
binning.

Sichuan is kept outside the baseline treated set because the 2022
hydro-thermal drought and power rationing are a non-EDWC shock on the
same system and would contaminate event-time coefficients at short
horizons. As a sensitivity, Sichuan is re-introduced with a
province-specific dummy for the 2022-08 to 2022-12 rationing window,
and the event-study coefficients are checked for stability.

Heterogeneity is reported along two cuts that the data support.
Cross-province heterogeneity is read directly from the cohort-specific
$\text{CATT}_{E_p, \ell}$ paths that feed the Sun--Abraham
aggregation, with attention to whether the four treated provinces
share a common sign and timing of response. A compute-tier cut uses
the CAICT zonghe suanli index to split treated-province months by
pre-EDWC compute capacity, following the provincial heterogeneity
documented in \citet{nietal-2024-co2emissionmitigationpathways}, so that
the reported effects can be read against the scale of pre-existing
compute infrastructure.

The CPC translation is re-computed under alternative settings of
$s_{\text{mxp}}(t)$, $\kappa_{\text{global}}$, $\lambda$, and the
PUE trajectory. The baseline Monte Carlo bands of §3.4.3 are
complemented by scenario runs that fix one parameter at an extreme
value while holding the others at their central values, so that the
marginal contribution of each assumption to the compute bridge is
visible.

A synthetic control exercise for Gansu is reported as a supplementary
robustness design rather than as a co-equal estimator. Following
\citet{abadieetal-2010-syntheticcontrolmethodsa} and
\citet{abadieetal-2015-comparativepoliticssynthetic}, a synthetic
Gansu is constructed from a donor pool of non-EDWC provinces that
excludes Sichuan, with pre-treatment predictors drawn from the
electricity and CO\textsubscript{2} panels, climate controls, and
pre-EDWC carbon intensity. Inference follows the in-space placebo
logic of \citet{abadieetal-2010-syntheticcontrolmethodsa}, comparing
pre- and post-treatment RMSPE against a placebo distribution.
Leave-one-out checks on donor weights are reported as a stability
diagnostic. The Sun--Abraham event study remains the main design;
agreement between the event-study path for Gansu and the synthetic
Gansu gap is read as convergent evidence, while disagreement is
reported honestly as a limitation of the identification rather than
as a rejection of either design.


% ------------------------------------------------------------
\subsection{Feasibility and Limitations}
% PURPOSE: Short, honest assessment.
%
% FEASIBLE:
% - CO2 panel ready; Gansu electricity verified
% - Sun--Abraham implementable with 4 treated provinces
% - CPC method documented and auditable
%
% LIMITATIONS:
% - Electricity coverage beyond Gansu in progress
% - No intraday data
% - Attribution risk: extra kWh may include non-DC load
% - CPC is intensive, not a measure of total AI compute
% - Treatment dates are working defaults
%
% FUTURE WORK (one sentence only):
% A possible extension would link estimated energy and compute
% effects to provincial economic outcomes such as fixed
% investment and ICT employment.
%
% Target: 0.5--1 page. FULLY DRAFTABLE.

The design described above is feasible with the data and methods that
are already in place. The provincial CO\textsubscript{2} panel built
from the Carbon Monitor China series of
\citet{liuetal-2020-carbonmonitornearrealtimea} covers every province
in the analysis window at a monthly frequency, so the emissions
outcome can be estimated from the first draft of the empirical
chapter. The Gansu electricity series has been rebuilt from the
official provincial bulletins and currently yields a seventy-two-month
monthly panel of total electricity consumption covering 2020-03
through 2026-02, with provenance tracked at the bulletin level. The
Sun--Abraham estimator is implementable on the four-province treated
set given the monthly calendar coverage and the staggered go-live
dates coded in §3.3.3, and the Carbon-per-Compute translation is
fully documented from the methodology note in the repository, so the
compute bridge can be rerun end-to-end by a third party.

The electricity analysis is planned for all four treated provinces.
Gansu is the first province for which the bulletin-scraping workflow
has been completed end to end, and it is used here to show that the
pipeline produces a province-level monthly series of sufficient
quality for the event study. The same workflow is being extended to
Ningxia, Inner Mongolia, and Guizhou one province at a time, as
recorded in the electricity data pipeline note; each extension reuses
the parsing logic and the provenance conventions established for
Gansu. Until the remaining three provinces are scraped, the
electricity outcome is reported in two stages: a single-province
event study for Gansu as a proof of concept, and a four-province
panel event study once the full panel is assembled. The
CO\textsubscript{2} outcome, which relies on the Carbon Monitor China
series rather than on bulletin scraping, is available for all four
treated provinces from the first draft of the empirical chapter. A second limitation is the
monthly frequency. No intraday load or price data are available, so
the design cannot speak to peak-shaving, curtailment reduction, or
wholesale-price effects; these are explicitly out of scope and were
dropped from the research plan for that reason.

An attribution risk remains even when the electricity panel is
complete. The estimated $\overline{\Delta\text{MWh}}_{p}$ captures the
full provincial load response in the months after go-live, not only
data-center demand, so part of the measured response may reflect
co-located industrial expansion or other EDWC-adjacent activity. The
translation in §3.4.3 addresses this by introducing the
data-center share parameter $\theta$ explicitly and reporting
compute-bridge outputs as bands rather than as point estimates. A
related limitation is that CPC is an intensive metric: it tracks
tCO\textsubscript{2} per exaflop-hour, not the total volume of AI
compute produced in a province, and does not by itself measure
national AI capacity. The coded go-live months are working defaults
drawn from policy announcements and public commissioning news; the
robustness exercises in §3.4.4 assess how far the main conclusions
move when those dates are shifted.

A possible extension, noted here but not pursued in the proposal,
would link the estimated energy and compute effects to provincial
economic outcomes such as fixed-asset investment and ICT employment.

\subsection{Research Structure}
% Research structure
%This paper mainly consists of six chapters. Chapter 1 is the introduction, mainly including problem statement...; Chapter 2 is the literature review, including...; Chapter 3 is the theoretical mechanism and research hypotheses... The following is the technical roadmap (the example below is simplified and needs to be more detailed):
%Introduction
%Literature Review
%Theoretical Mechanism and Research Hypotheses
%Research Design
%Empirical Results
%Conclusion and Recommendations

The thesis will be organized into nine chapters that follow the
logic of the two research questions.

\begin{itemize}
  \item \textbf{Chapter 1} will motivate the topic, situate EDWC at
        the intersection of China's compute build-out and its
        dual-carbon commitments, and state the two research questions.
  \item \textbf{Chapter 2} will review the institutional background
        of EDWC, the relationship between data-center operation and
        electricity demand, and the Chinese policy framework for
        green compute.
  \item \textbf{Chapter 3} will document the monthly province-level
        panel: the Carbon Monitor China CO\textsubscript{2} series,
        the scraped provincial electricity bulletins, the PUE and
        AI-equivalent efficiency series, and the coded EDWC go-live
        dates.
  \item \textbf{Chapter 4} will set out the identification strategy
        and the Sun--Abraham event-study specification for RQ1, with
        the synthetic-control exercise for Gansu as a supplementary
        check.
  \item \textbf{Chapter 5} will report the event-study results for
        electricity and CO\textsubscript{2}, together with
        cross-province and compute-tier heterogeneity.
  \item \textbf{Chapter 6} will introduce the Carbon-per-Compute
        framework and translate the reduced-form estimates into the
        three compute views defined in §3.4.3, with Monte Carlo
        uncertainty bands.
  \item \textbf{Chapter 7} will collect robustness exercises on
        alternative treatment timings, donor pools, and CPC
        parameter settings.
  \item \textbf{Chapter 8} will discuss the implications of the
        results for EDWC strategy and present the exploratory
        extension to provincial economic outcomes; that extension
        will be framed as descriptive and program-adjacent rather
        than as a third causal question.
  \item \textbf{Chapter 9} will conclude with the policy implications
        and a short agenda for future work on hourly data, prices,
        and curtailment.
\end{itemize}

Technical appendices will collect the full formulas for GF/W, PUE,
and CPC, the treatment-date source tables, and the Monte Carlo
details.


%============================================================
% SECTION 4 --- CONTRIBUTIONS AND INNOVATIONS
% ============================================================
% Target length: 1--1.5 pages

\section{Contributions and Innovations}

% ------------------------------------------------------------
\subsection{Academic Contributions}
% TWO MAIN POINTS:
%
% 1. First study to apply staggered-adoption event-study
%    (Sun--Abraham) to EDWC go-lives, estimating province-level
%    causal effects on electricity consumption and CO2 emissions.
%    Fills the gap between simulation-based EDWC studies and the
%    need for credible causal evidence.
%
% 2. Construction of a reproducible Carbon-per-Compute (CPC) index
%    combining provincial carbon intensity, PUE, and an AI-equivalent
%    efficiency series with China-specific export-control adjustments.
%    Provides a transparent energy-to-compute bridge applicable
%    beyond this thesis.
%
% DO NOT OVERCLAIM: call CPC a "transparent proxy," not a
% definitive metric.

The thesis proposes two academic contributions, one on the design
side and one on the measurement side.

The first contribution is on the design of EDWC evaluation. Existing
province-level assessments of EDWC rely on engineering simulations or
forward-looking scenarios, as in the migration-and-routing simulation
of \citet{zhangetal-2025-decarbonizingdatacenters} and the net-zero
scenario assessment of \citet{zhangetal-2025-easterndatawestern};
these establish the magnitudes that an EDWC reallocation could
plausibly deliver, but they do not estimate what the program has
delivered in the provinces that went live. Applying the cohort-interacted
event-study estimator of
\citet{sunabraham-2021-estimatingdynamictreatmenta} to the four
treated provinces, with the
\citet{callawaysantanna-2021-differenceindifferencesmultipletimea}
and \citet{borusyaketal-2024-revisitingeventstudydesigns} estimators
reported as robustness, moves the EDWC literature from simulation to
ex post reduced-form estimates of the monthly electricity and
CO\textsubscript{2} response, with pre-trend diagnostics and
cohort-specific paths rather than a single pooled coefficient.

The second contribution is on measurement. The thesis constructs a
transparent Carbon-per-Compute index that links provincial carbon
intensity, facility power usage effectiveness, and an AI-equivalent
hardware efficiency series into a single tCO\textsubscript{2} per
exaflop-hour metric, following the carbon-aware computing framework
of \citet{guptaetal-2022-chasingcarbonelusive}. The specific
innovation is the China-layer construction of
$\text{GF/W}_{\text{AI,China}}(t)$: the Green500 FP64 frontier is
bridged to AI-like mixed precision through the HPL-MxP speedup, a
utilization realism factor is applied, and a time-varying
interconnect haircut captures the A800 and H800 export-control
regimes. Every input is observable or benchmarkable, so the metric
can be recomputed under alternative assumptions by other researchers.

% ------------------------------------------------------------
\subsection{Policy Contributions}
% TWO MAIN POINTS:
%
% 1. Province-level evidence on EDWC energy and emissions effects
%    that can inform compute-hub siting decisions and regional
%    energy planning under dual-carbon constraints.
%
% 2. A CPC-based provincial ranking framework that could support
%    green compute procurement, data-center carbon accounting,
%    and compute-electricity synergy policy.
%
% Keep compact. Do not oversell.

The proposal is designed to produce two inputs to compute and
energy policy in China, at the provincial level that EDWC actually
operates on.

The first policy contribution is province-level reduced-form
evidence on how EDWC has moved electricity consumption and
CO\textsubscript{2} emissions in the four treated provinces after
go-live. Policy assessments of EDWC to date have relied on
forward-looking scenarios
\citep{zhengetal-2020-mitigatingcurtailmentcarbon,
zhangetal-2025-easterndatawestern}, which set an expected envelope
but leave the realized provincial response unmeasured. Event-study
estimates with pre-trend diagnostics and cohort-specific paths
report what happened in the provinces that went live first,
providing a direct input to compute-hub siting, provincial grid
planning, and the pacing of further EDWC expansion under the
dual-carbon targets.

The second policy contribution is the Carbon-per-Compute framework
as a provincial ranking and accounting device. Because CPC is built
from observable inputs, it can be updated as the grid mix, facility
efficiency, and hardware stock evolve, and its three channels (grid,
site, hardware) are kept separate by construction so that policy
attention can be directed at the component that is actually moving.
The framework is usable for green-compute procurement, data-center
carbon accounting under the NDRC and MIIT green-data-center programs,
and for tracking whether the compute-electricity synergy envisaged
in EDWC is being delivered in intensity terms.


% ============================================================
% SECTION 5 --- WORK ALREADY COMPLETED AND CURRENT PROGRESS
% ============================================================
% Target length: 1--1.5 pages
% Purpose: show the committee that real work has been done.
%
% CONTENT: describe completed items (CO2 panel, Gansu electricity,
% treatment-date archive, CPC method documentation, CAICT mapping)
% and remaining items (non-Gansu electricity, weather controls,
% estimation, literature completion).

\section{Work Already Completed and Current Progress}

Real work underlies the proposal. The core panel infrastructure, the
treatment-date archive, and the Carbon-per-Compute method note are
already in place; the main remaining task is extending the
bulletin-scraping pipeline from Gansu to the other three treated
provinces. Table~\ref{tab:progress} summarizes the status of each
work stream.

\begin{table}[h]
\centering
\small
\begin{tabular}{p{4.5cm} p{7.5cm} p{2cm}}
\toprule
\textbf{Work stream} & \textbf{Current status} & \textbf{State} \\
\midrule
Provincial CO\textsubscript{2} panel &
Carbon Monitor China series assembled at monthly frequency for 31
provinces, 2019-01 to 2025-11; ready for the event study. & Done \\
\addlinespace
Gansu electricity series &
Provincial bulletins scraped, parsed, and validated; seventy-two
monthly observations with bulletin-level provenance. & Done \\
\addlinespace
Electricity series for Ningxia, Inner Mongolia, Guizhou &
Scrapers being adapted province by province, reusing the Gansu
parsing logic and provenance conventions. & In progress \\
\addlinespace
EDWC treatment-date archive &
Four go-live months coded from policy announcements and public
commissioning news; sources tracked in
\texttt{01\_notes/edwc\_treatment\_dates.md}. & Done \\
\addlinespace
CPC method documentation &
\texttt{01\_notes/AI\_equivalent\_GFLOPS\_per\_W.md} specifies the
Green500 frontier, HPL-MxP speedup, $\kappa_{\text{global}}$, and
the interconnect haircut construction. & Done \\
\addlinespace
CAICT composite compute index &
Provincial \emph{zonghe suanli zhishu} series mapped for the
heterogeneity cuts; CAICT development index tracked separately. & Done \\
\addlinespace
Climate and holiday controls &
Data sources identified; covariates to be finalized in the empirical
chapter. & In progress \\
\addlinespace
Event-study estimation &
Specification finalized in §3.4.2; estimation pending completion of
the four-province electricity panel. & Pending \\
\addlinespace
Literature review &
\texttt{02\_literature/literature\_master.md} and the summary notes
cover the five streams; additional screening will accompany the
empirical chapter. & Ongoing \\
\bottomrule
\end{tabular}
\caption{Status of the main work streams supporting the proposal.}
\label{tab:progress}
\end{table}


% ============================================================
% SECTION 6 --- RESEARCH PLAN AND TIMELINE
% ============================================================
% Target length: 0.5--1 page
% Purpose: show a phased completion plan.
%
% CONTENT: 5 phases --- data collection, baseline estimation,
% CPC construction, writing, revision/defense.

\section{Research Plan and Timeline}

[TO BE DRAFTED]


% ============================================================
% REFERENCES
% ============================================================

\bibliographystyle{apalike}
\bibliography{references}


\end{document}